
\section{OA-05: Conditional filtering can distinguish non-orthogonal states}
\source{Week 6 Tutorial, Question 2; released Week 6 tutorial solutions.}
\subsection*{Step 1: measure the input overlap}
The inputs and successful Kraus operator are
\[
\ket{\psi_\pm}=\frac{2\ket0\pm\ket1}{\sqrt5},\qquad
K_s=\tfrac12\ket0\bra0+\ket1\bra1 .
\]
Both input states are normalised, but
\[
\boxed{\langle\psi_+|\psi_-\rangle
=\frac{4-1}{5}=\frac35\ne0}.
\]
Thus no \emph{deterministic} measurement can perfectly
distinguish them.

\subsection*{Step 2: apply the success operator before normalising}
\[
K_s\ket{\psi_\pm}=\frac{\ket0\pm\ket1}{\sqrt5}.
\]
The trace/weight of this unnormalised branch is
\[
\boxed{p_s=\|K_s\ket{\psi_\pm}\|^2
=\frac15+\frac15=\frac25},
\]
the same for either input. Conditional on success,
\[
\boxed{\ket{\phi_\pm}=
\frac{K_s\ket{\psi_\pm}}{\sqrt{p_s}}
=\frac{\ket0\pm\ket1}{\sqrt2}},
\quad
\langle\phi_+|\phi_-\rangle=0 .
\]
A subsequent measurement in the \(\{\ket+,\ket-\}\) basis
identifies which input was sent \emph{when the filter succeeds}.

\subsection*{Step 3: explain the failure outcome}
The other Kraus operator is
\(K_f=(\sqrt3/2)\ket0\bra0\).
Check completeness:
\[
K_s^\dagger K_s+K_f^\dagger K_f
=\begin{pmatrix}1/4&0\\0&1\end{pmatrix}
+\begin{pmatrix}3/4&0\\0&0\end{pmatrix}=I.
\]
Each failed branch is
\[
K_f\ket{\psi_\pm}=\sqrt{\frac35}\ket0,
\quad p_f=\frac35,\quad
\ket{\phi_f}=\ket0.
\]
Failure produces the \emph{same} state for both signs
and reveals no information distinguishing the inputs.
\checkmethod{The output states become orthogonal only in the
successful subensemble. It is the non-unit success probability
that prevents violation of deterministic distinguishability.}
